% Generated draft: requires mathematical and source review.
\documentclass[11pt]{article}
\usepackage[margin=25mm]{geometry}
\usepackage{fontspec}
\setmainfont{Latin Modern Roman}
\usepackage{amsmath,amssymb,mathtools,unicode-math}
\setmathfont{Latin Modern Math}
\usepackage{enumitem,needspace,xurl,hyperref}
\setlength{\parindent}{0pt}
\setlength{\parskip}{5pt}
\setlength{\emergencystretch}{3em}
\newcommand{\N}{\mathbb{N}}
\newcommand{\Z}{\mathbb{Z}}
\newcommand{\Q}{\mathbb{Q}}
\newcommand{\R}{\mathbb{R}}
\newcommand{\C}{\mathbb{C}}
\newcommand{\F}{\mathbb{F}}
\newcommand{\A}{\mathbb{A}}
\newcommand{\PP}{\mathbb{P}}
\newcommand{\E}{\mathbb{E}}
\newcommand{\eps}{\varepsilon}
\newcommand{\dd}{\,\mathrm{d}}
\begin{document}
\section*{UnsolvedMath ID: 30005543}
\textbf{Source problem code:} catalogue entry
\subsection{Definitions and mathematical statement}
Let \(S^1=\{\omega\in\R^2:|\omega|=1\}\). For explicit notation, take \(\sigma\) to be arclength measure and use the Fourier transform convention
\[
\widehat{f\sigma}(x)=\int_{S^1}f(\omega)e^{-i x\cdot\omega}\,\dd\sigma(\omega),\qquad x\in\R^2.
\]
These measure and Fourier conventions are stated here because the supplied record does not specify them; the question of whether constants attain the optimal norm is unaffected by the usual changes of these conventions.

For a scalar field \(\mathbb K\in\{\R,\C\}\), define the optimal Tomas--Stein extension constant by
\[
C_{\mathbb K}=\sup_{0\ne f\in L^2(S^1,\sigma;\mathbb K)}
\frac{\|\widehat{f\sigma}\|_{L^6(\R^2)}}{\|f\|_{L^2(S^1,\sigma)}}.
\]
A nonzero function is an extremizer if its quotient equals \(C_{\mathbb K}\).

\textbf{Question.} For the scalar field intended by the source, do nonzero constant functions attain this optimal constant? Equivalently, is
\[
C_{\mathbb K}=\frac{\|\widehat{\sigma}\|_{L^6(\R^2)}}{\sqrt{2\pi}},
\]
or, equivalently, does every \(f\in L^2(S^1,\sigma;\mathbb K)\) satisfy
\[
\|\widehat{f\sigma}\|_{L^6(\R^2)}
\leq \frac{\|\widehat{\sigma}\|_{L^6(\R^2)}}{\sqrt{2\pi}}
\|f\|_{L^2(S^1,\sigma)}?
\]
The source does not explicitly specify real-valued versus complex-valued functions. The parameter \(\mathbb K\) exposes that unresolved choice; this draft does not assert an equivalence between the two formulations. The question concerns attainment by constants, without asserting uniqueness or classifying all extremizers.

\textbf{Qualifications from the supplied original excerpt.} A function \(g\) on \((S^1)^3\) is called antipodal when it is even in each argument separately: for each \(j\in\{1,2,3\}\), replacing \(\omega_j\) by \(-\omega_j\) leaves \(g\) unchanged almost everywhere. The excerpt discusses a reduction of the conjectured sharp inequality to positive semidefiniteness of a quadratic form on antipodal functions in \(L^2((S^1)^3,\sigma^{\otimes3})\). Its corrupted transcription suggests the formal expression
\[
\begin{aligned}
Q(g)={}&\int_{(S^1)^6}
\delta(\omega_1+\omega_2+\omega_3+\omega_4+\omega_5+\omega_6)
\bigl(|\omega_1+\omega_2+\omega_3|^2-1\bigr)\\
&\quad\times\bigl(|g(\omega_1,\omega_2,\omega_3)|^2
-g(\omega_1,\omega_2,\omega_3)g(\omega_4,\omega_5,\omega_6)\bigr)
\prod_{j=1}^{6}\dd\sigma(\omega_j).
\end{aligned}
\]
Here the \(\delta\) notation is retained from the excerpt as a formal Dirac constraint. The supplied fragment does not provide a rigorous interpretation or domain for this singular expression, and it supplies no complex conjugate in the mixed product. Consequently, this display is a transcription aid, not a complete definition of a quadratic form on complex \(L^2\), and no precise equivalence with the main question is asserted here.

The excerpt attributes to Barker, Thiele and Zorin-Kranich numerical verification of positive semidefiniteness on antipodal functions with Fourier modes up to degree \(120\), and reports concentration of eigenfunctions associated with small eigenvalues in that finite-dimensional computation. The meaning of the degree cutoff is not specified. It also reports positive semidefiniteness on a certain unspecified subspace containing functions with localized support, with a corollary giving the conjectured sharp inequality for functions supported near a pair \(\{p,-p\}\), \(p\in S^1\). The printed neighborhood size contains \(6/80\) and a detached square-root sign; neither its exact radius nor its distance convention can be recovered reliably from this record.

The supplied literature assessment, dated \(2026\text{-}08\text{-}22\), attributes the localized sharp inequality to Becker and labels the overall problem partially solved, with the global constant-extremizer conjecture remaining open. These are supplied dated curation claims, not independently verified mathematical or literature conclusions.
\subsection{Short English statement}
Does a nonzero constant function maximize the ratio of the \(L^6(\R^2)\) norm of the circle extension operator to the \(L^2(S^1,\sigma)\) norm of its input? The supplied record reports a localized partial result, while leaving the global question open; the scalar field and several details of the auxiliary quadratic-form formulation require review.
\subsection{Sources}
\begin{itemize}
\item Supplied local registry record: \texttt{data/unsolvedmath/problems.json}, permanent UnsolvedMath ID \(30005543\), exact source problem code \texttt{catalogue entry}, title \emph{Tomas--Stein Extremizers on the Circle}. The record gives proposed year \(2023\) and names no proposer.
\item Supplied report citation: \emph{Incidence Problems in Harmonic Analysis, Geometric Measure Theory, and Ergodic Theory} (2023), Oberwolfach Reports, pp.~1397--1452, DOI \texttt{10.4171/owr/2023/25}. \url{https://doi.org/10.4171/owr/2023/25}. The supplied background describes these reports as documenting workshops and their open-problem discussions.
\item Supplied literature source, labeled ``OWR 2023 25'': \url{https://publications.mfo.de/bitstream/handle/mfo/4055/OWR_2023_25.pdf?sequence=1}.
\item Supplied literature source, labeled ``Becker'': \url{https://arxiv.org/abs/2304.02345}. The supplied assessment attributes a sharp inequality for functions localized near an antipodal pair to Becker; no full bibliographic description is supplied.
\item Supplied literature source, labeled ``Barker et al'': \url{https://arxiv.org/abs/2002.05118}. The original excerpt attributes the numerical computations to Barker, Thiele and Zorin-Kranich and labels that reference [1]; it also refers to recent work [2], without supplying the numbered bibliography in the excerpt.
\item Literature checked at: \texttt{2026-08-22}. The supplied dated curation labels the problem \texttt{partially\_solved} and states that the global constant-extremizer conjecture remains open. This assessment and its cited claims have not been independently verified here; the supplied background explicitly requests independent verification of time-sensitive status and claims.
\item The registry labels the statement \texttt{corrected\_verified} and reports that it was checked against the official Oberwolfach report. This is supplied verification metadata, not a verification performed for this draft.
\end{itemize}
\end{document}
